\chapter{The Type-Law Architecture}
\label{ch:architecture}

This chapter describes the three-layer type system of \texttt{wasm4pm-compat}
in precise detail: the evidence lifecycle (state tokens and the one-way door),
the witness marker system (uninhabited authority labels and the sealed
\texttt{Witness} trait), the const-generic law machinery
(\texttt{Assert}/\texttt{IsTrue}/\texttt{Require} and \texttt{Between01}), the
non-forgeability mechanism (\texttt{WfNetConst} and sealed constructors), and
the loss-policy chain (\texttt{LossPolicy}, \texttt{ProjectionName}, and
\texttt{LossReport}).

Every invariant described here is enforced by the Rust type system at compile
time. No invariant is a runtime assertion, a documentation comment, or a
convention agreed by convention alone.

% -----------------------------------------------------------------------------
\section{The Evidence Lifecycle}
\label{sec:arch-lifecycle}
% -----------------------------------------------------------------------------

\subsection{The One-Way Door}

The central structural invariant of \texttt{wasm4pm-compat} is a typed,
one-way process-evidence lifecycle. Raw, untrusted input cannot be used as
admitted evidence; admitted evidence cannot be un-admitted; refused evidence
cannot be rehabilitated. The lifecycle is directed and acyclic, and it is
enforced entirely by the Rust type system.

\begin{definition}[Evidence Lifecycle]
  \label{def:lifecycle}
  The \emph{evidence lifecycle} is the directed state machine over the set of
  stage tokens
  $\mathcal{S} = \{\texttt{Raw}, \texttt{Parsed}, \texttt{Admitted},
  \texttt{Refused}, \texttt{Projected}, \texttt{Exportable}, \texttt{Receipted}\}$
  with initial state \texttt{Raw} and terminal states
  $\{\texttt{Refused}, \texttt{Receipted}\}$.
\end{definition}

The lawful transitions are:

\begin{equation}
  \label{eq:lifecycle}
  \texttt{Raw}
  \xrightarrow{\texttt{into\_parsed}}
  \texttt{Parsed}
  \xrightarrow{\texttt{Admit::admit}}
  \texttt{Admitted}
  \xrightarrow{\quad}
  \begin{cases}
    \texttt{Projected}   & \xrightarrow{\texttt{into\_receipted}} \texttt{Receipted} \\
    \texttt{Exportable}  & \xrightarrow{\texttt{into\_receipted}} \texttt{Receipted} \\
    \texttt{Receipted}   & \text{(terminal)}
  \end{cases}
\end{equation}

The refuse path is available at every stage before admission:

\begin{equation}
  \label{eq:refuse}
  \texttt{Raw} \xrightarrow{\texttt{refuse}} \texttt{Refused}
  \qquad
  \texttt{Parsed} \xrightarrow{\texttt{into\_refused}} \texttt{Refused}
\end{equation}

\texttt{Refused} is terminal. No transition out of \texttt{Refused} exists.

\subsection{PhantomData Encoding}

The stage tokens are zero-sized empty enums used exclusively as \texttt{PhantomData}
type parameters. They contribute zero bytes to the memory layout of
\Evidence{T}{State}{W}. The full lifecycle enforcement is erased at compile time.

\begin{lstlisting}[language=Rust, caption={The \Evidence{}{}{} carrier with
  zero-sized stage and witness tags. \texttt{State} and \texttt{W} are pure
  compile-time phantoms.}, label={lst:evidence-struct}]
pub struct Evidence<T, State: EvidenceState, W> {
    pub value: T,
    pub state:   PhantomData<State>,
    pub witness: PhantomData<W>,
}
\end{lstlisting}

Because \texttt{State} is a type parameter and not a runtime value,
\texttt{Evidence<T, Raw, W>} and \texttt{Evidence<T, Admitted, W>} are
\emph{different types}. They share no subtype relationship, no
\texttt{From}/\texttt{Into} implementation, and no common constructor reachable
from outside the crate. The \texttt{EvidenceState} bound is a sealed trait
implemented only for the seven canonical stage tokens; a downstream crate
cannot introduce a new stage.

\subsection{Compile-Time Proof That Cross-State Substitution Is a Type Error}

The type-level safety property is:

\begin{equation}
  \label{eq:no-subtype}
  \forall T, S_1, S_2, W.\quad S_1 \neq S_2
  \implies
  \Evidence{T}{S_1}{W} \not\leq \Evidence{T}{S_2}{W}
\end{equation}

This is not merely a claim about runtime behavior; it is a consequence of
Rust's nominal type system. Two instantiations of the same generic type with
different type arguments are different types with no implicit coercion between
them.

The trybuild compile-fail fixture for this property is:

\begin{lstlisting}[language=Rust, caption={Compile-fail fixture: assigning
  \texttt{Raw} evidence to an \texttt{Admitted} binding is a type error.
  The error fires for the named lifecycle law, not accidentally.},
  label={lst:cf-raw-as-admitted}]
// tests/ui/compile_fail/raw_used_as_admitted.rs
use wasm4pm_compat::evidence::Evidence;
use wasm4pm_compat::state::{Raw, Admitted};
use wasm4pm_compat::witness::Ocel20;

fn require_admitted(_: Evidence<u32, Admitted, Ocel20>) {}

fn main() {
    let raw: Evidence<u32, Raw, Ocel20> = Evidence::raw(42u32);
    require_admitted(raw);
    //               ^^^
    // error[E0308]: mismatched types
    //   expected `Evidence<u32, Admitted, Ocel20>`
    //      found `Evidence<u32, Raw, Ocel20>`
}
\end{lstlisting}

The \texttt{Admitted} constructor is \texttt{pub(crate)}. It is called only
from within the \texttt{admission} module, exclusively by implementors of
\texttt{Admit::admit}. There is no public free function, no \texttt{From}
impl, and no unsafe transmute path that produces \texttt{Evidence<T, Admitted,
W>} from outside the crate.

\subsection{Full Lifecycle Walk-Through}

Listing~\ref{lst:full-lifecycle} shows the complete lawful path from \texttt{Raw}
to \texttt{Receipted} as it appears in the crate's test suite.

\begin{lstlisting}[language=Rust, caption={Full lifecycle: \texttt{Raw}
  $\to$ \texttt{Admitted} $\to$ \texttt{Projected} $\to$ \texttt{Receipted}.
  Each transition is a distinct method on a distinct type. None can be
  skipped without a compiler error.}, label={lst:full-lifecycle}]
use wasm4pm_compat::admission::Admission;
use wasm4pm_compat::evidence::Evidence;
use wasm4pm_compat::state::{Raw, Admitted, Projected, Receipted};
use wasm4pm_compat::witness::Ocel20;

// Step 1: Raw evidence --- freely constructible.
let raw: Evidence<u32, Raw, Ocel20> = Evidence::raw(1u32);

// Step 2: Parse (format-acceptance, not yet structurally judged).
let parsed = raw.into_parsed();

// Step 3: Admission --- only via Admit::admit(). Simulated here via
// Admission::new(), which is what an Admit impl calls internally.
let admitted: Evidence<u32, Admitted, Ocel20> =
    Admission::new(1u32).into_evidence();

// Step 4: Named lossy projection.
let projected: Evidence<u32, Projected, Ocel20> =
    admitted.into_projected();

// Step 5: Receipt --- terminal stage.
let receipted: Evidence<u32, Receipted, Ocel20> =
    projected.into_receipted();

assert_eq!(receipted.value, 1u32);
\end{lstlisting}

% -----------------------------------------------------------------------------
\section{The Witness Marker System}
\label{sec:arch-witness}
% -----------------------------------------------------------------------------

\subsection{Authority Labels as Uninhabited Types}

A \emph{witness marker} is an empty enum --- an uninhabited type. It names
a specific standard, paper, API grammar, Rust law, or internal bridge that a
piece of evidence is being admitted, projected, or graduated \emph{against}.
Witnesses carry no data and do no work; their sole function is to thread a
named authority through the type signature.

\begin{definition}[Witness Marker]
  \label{def:witness}
  A \emph{witness marker} is an uninhabited type $W$ implementing the sealed
  \texttt{Witness} trait. It is associated with a machine-facing key
  (\texttt{KEY}), a human-facing title (\texttt{TITLE}), a publication year
  (\texttt{YEAR}), and a family classification (\texttt{FAMILY}).
\end{definition}

\subsection{The Sealed Witness Pattern}

The \texttt{Witness} trait is publicly visible but cannot be implemented
outside the crate. This sealing is achieved by the macro-generated
\texttt{pub enum WitnessMarker \{\}} pattern: the macro declares an empty
public enum (which is uninhabited and therefore never constructible) and
immediately provides the trait implementation. Because the enum itself lives in
the crate, external code cannot create a new instance and cannot implement
\texttt{Witness} for a type it controls.

\begin{lstlisting}[language=Rust, caption={The \texttt{Witness} trait and
  the \texttt{witness\_marker!} macro that seals it. Every marker is an
  uninhabited empty enum.}, label={lst:witness-sealed}]
pub trait Witness {
    const KEY:    &'static str;
    const FAMILY: WitnessFamily;
    const TITLE:  &'static str;
    const YEAR:   Option<u16>;
}

// Macro-generated: each invocation produces an uninhabited type + impl.
macro_rules! witness_marker {
    ($(#[$meta:meta])* $name:ident,
     $key:literal, $family:expr, $title:literal, $year:expr) => {
        $(#[$meta])*
        #[derive(Debug, Clone, Copy, PartialEq, Eq, Hash)]
        pub enum $name {}            // uninhabited --- never constructible

        impl Witness for $name {
            const KEY:    &'static str = $key;
            const FAMILY: WitnessFamily = $family;
            const TITLE:  &'static str = $title;
            const YEAR:   Option<u16>  = $year;
        }
    };
}

// Representative witnesses (27+ defined in witness.rs):
witness_marker!(Ocel20,
    "ocel-2.0", WitnessFamily::Standard, "OCEL 2.0", Some(2023));
witness_marker!(Xes1849,
    "xes-1849-2016", WitnessFamily::Standard,
    "XES (IEEE 1849-2016)", Some(2016));
witness_marker!(WfNetSoundnessPaper,
    "wfnet-soundness-paper", WitnessFamily::Paper,
    "WF-net soundness (van der Aalst 1998)", Some(1998));
witness_marker!(InductiveMiner,
    "inductive-miner", WitnessFamily::Paper,
    "Inductive Miner (Leemans et al. 2013)", Some(2013));
\end{lstlisting}

The crate declares 38 witness markers in \texttt{src/witness.rs}, spanning
five families: \texttt{Standard} (OCEL~2.0, XES, XES sub-extensions),
\texttt{Paper} (POWL, WF-net soundness, OC-Petri nets, OCPQ, Inductive Miner,
Alpha Miner, Declare, Log Skeleton, YAWL, alignment, process cubes, and more),
\texttt{ApiGrammar} (pm4py, pmax), \texttt{RustLaw} (typestate law), and
\texttt{InternalBridge} (wasm4pm graduation bridge).

\subsection{Why \texttt{Admission<T, Ocel20>} and \texttt{Admission<T, Xes1849>} Are Genuinely Distinct}

Because \texttt{Ocel20} and \texttt{Xes1849} are different types (even though
both are empty enums), \texttt{Admission<T, Ocel20>} and
\texttt{Admission<T, Xes1849>} are different instantiations of the generic
\texttt{Admission<T, W>} struct --- with no common supertype and no implicit
conversion between them.

A function written to process OCEL~2.0 evidence:

\begin{lstlisting}[language=Rust, label={lst:witness-distinct}]
fn process_ocel(e: Admission<OcelLog, Ocel20>) { /* ... */ }
\end{lstlisting}

\noindent cannot be called with an XES admission:

\begin{lstlisting}[language=Rust]
let xes_admission: Admission<XesLog, Xes1849> = /* ... */;
process_ocel(xes_admission);
// error[E0308]: mismatched types
//   expected `Admission<OcelLog, Ocel20>`
//      found `Admission<XesLog, Xes1849>`
\end{lstlisting}

This is not a runtime type check; it is a compile-time identity check on the
type parameter. The witness type is erased to zero bytes at runtime --- the
enforcement cost is entirely compile-time.

% -----------------------------------------------------------------------------
\section{Const-Generic Law Machinery}
\label{sec:arch-constlaw}
% -----------------------------------------------------------------------------

\subsection{The Assert/IsTrue/Require Pattern}

The \texttt{Assert}/\texttt{IsTrue}/\texttt{Require} pattern is the mechanism
by which boolean const-generic expressions become compile-time proof
obligations. It requires the nightly \texttt{generic\_const\_exprs} feature.

\begin{lstlisting}[language=Rust, caption={The compile-time bounds machinery.
  \texttt{Assert<true>} is inhabited; \texttt{Assert<false>} is not.
  \texttt{Require<\{EXPR\}>: IsTrue} becomes a compile error when \texttt{EXPR}
  evaluates to \texttt{false}.}, label={lst:assert-istrue}]
pub struct Assert<const OK: bool>;

pub trait IsTrue {}
impl IsTrue for Assert<true> {}      // only the true case is inhabited

pub type Require<const OK: bool> = Assert<OK>;
\end{lstlisting}

A where-bound of the form \texttt{Require<\{ EXPR \}>: IsTrue} compiles
successfully if and only if \texttt{EXPR} evaluates to \texttt{true} at
compile time. When \texttt{EXPR} is \texttt{false}, the compiler reports:

\begin{verbatim}
error[E0277]: the trait bound `Assert<false>: IsTrue` is not satisfied
\end{verbatim}

This is the mechanism behind every quantitative invariant in the crate:
\texttt{ConditionCell<BITS>}, \texttt{Between01<NUM, DEN>},
\texttt{TypedLoopNode<ARITY>}, and the \texttt{WfNetConst<SOUNDNESS>}
non-forgeability gate all reduce to \texttt{Require}/\texttt{IsTrue} bounds.

\subsection{Between01 for Metric Bounds}

\texttt{Between01<NUM, DEN>} encodes the invariant that the rational number
$\texttt{NUM}/\texttt{DEN}$ lies in $[0, 1]$ at the type level, with no
runtime component.

\begin{lstlisting}[language=Rust, caption={\texttt{Between01<NUM, DEN>}: a
  rational metric in $[0,1]$ enforced by two \texttt{Require}/\texttt{IsTrue}
  where-bounds.}, label={lst:between01}]
pub struct Between01<const NUM: u64, const DEN: u64>
where
    Require<{ DEN > 0 }>:   IsTrue,  // denominator non-zero
    Require<{ NUM <= DEN }>: IsTrue,  // numerator at most denominator
{
    _private: (),
}
\end{lstlisting}

The two where-bounds are independent:
\begin{itemize}
  \item \texttt{Between01<0, 0>} does not compile: $\texttt{DEN} = 0$ violates
    the non-zero-denominator law.
  \item \texttt{Between01<2, 1>} does not compile: $2 > 1$ violates the upper
    bound law.
  \item \texttt{Between01<3, 4>} compiles: $3/4 \in [0,1]$ and $4 > 0$.
\end{itemize}

\begin{theorem}[Between01 Soundness]
  \label{thm:between01}
  For any $n, d \in \mathbb{N}$, the type $\texttt{Between01}\langle n, d \rangle$
  is well-formed if and only if $d > 0$ and $n \leq d$.
\end{theorem}

\begin{proof}
  The type definition carries two where-bounds:
  $\RequireBound{d > 0} : \IsTrue$ and $\RequireBound{n \leq d} : \IsTrue$.
  By the \texttt{Assert}/\texttt{IsTrue} machinery, each bound reduces to
  $\texttt{Assert}\langle b \rangle : \IsTrue$ where $b$ is the boolean value
  of the bracketed expression evaluated at compile time. The trait
  \texttt{IsTrue} is implemented only for $\texttt{Assert}\langle\texttt{true}\rangle$.
  Therefore, well-formedness holds if and only if both $d > 0$ and $n \leq d$
  evaluate to \texttt{true}, i.e., if and only if $d > 0$ and $n \leq d$.
  \qed
\end{proof}

\texttt{Between01} is used by \texttt{Metric<KIND, NUM, DEN>} in
\texttt{src/conformance.rs} to ensure that conformance quality metrics
(fitness, precision, F1, generalization, simplicity) are provably in-range at
the type level. A \texttt{Metric<Fitness, 5, 3>} --- claiming a fitness of
$5/3 > 1$ --- does not compile.

\subsection{ConditionCell and the Need-9 Law}

\texttt{ConditionCell<BITS>} encodes the \emph{Need9 means split} law from the
Blue River Dam covenant: a single condition cell may hold at most 8 primary
bits. When you need 9 primary condition bits, the law requires splitting into
two cells.

\begin{lstlisting}[language=Rust, caption={\texttt{ConditionCell<BITS>}: the
  Need9 law as a type constraint.}, label={lst:conditioncell}]
pub struct ConditionCell<const BITS: usize>
where
    Require<{ BITS <= 8 }>: IsTrue,
{
    _private: (),
}
\end{lstlisting}

\texttt{ConditionCell<8>} compiles. \texttt{ConditionCell<9>} does not:

\begin{verbatim}
error[E0277]: the trait bound `Assert<false>: IsTrue` is not satisfied
  --> src/law.rs
   |
   |     Require<{ BITS <= 8 }>: IsTrue,
   |                             ^^^^^^ the trait `IsTrue` is not
   |                                    implemented for `Assert<false>`
\end{verbatim}

The error message is accompanied by a developer note in the source directing
the reader to the Blue River Dam covenant and the split remedy.

% -----------------------------------------------------------------------------
\section{Non-Forgeability via Sealed Constructors}
\label{sec:arch-nonforgeability}
% -----------------------------------------------------------------------------

\subsection{The WfNetConst Private Seal}

\texttt{WfNetConst<SOUNDNESS>} in \texttt{src/petri.rs} is the typed
representation of a workflow net whose soundness state is tracked as a const
generic parameter drawn from the \texttt{SoundnessState} enum:
$\{\texttt{Unknown}, \texttt{Claimed}, \texttt{Witnessed}\}$.

The non-forgeability of the \texttt{Witnessed} state is enforced by a private
constructor seal. Only the \texttt{witness\_soundness} function inside the
\texttt{petri} module may mint a \texttt{WfNetConst<\{Witnessed\}>}; it
requires a \texttt{SoundnessProof} token as proof that the engine has verified
soundness. External code has no path to \texttt{WfNetConst<\{Witnessed\}>}
that does not pass through this gate.

\begin{lstlisting}[language=Rust, caption={Sealed constructor for
  \texttt{WfNetConst<\{Witnessed\}>}. The \texttt{SoundnessProof} token
  is only producible by the \texttt{wasm4pm} engine bridge. External
  code cannot construct it.}, label={lst:wfnet-seal}]
pub struct WfNetConst<const SOUNDNESS: SoundnessState> {
    // fields ...
    _seal: SoundnessSeal<SOUNDNESS>,
}

// Private seal type --- not exported.
struct SoundnessSeal<const S: SoundnessState>;

// The only path to WfNetConst<{Witnessed}>:
pub fn witness_soundness(
    net:   WfNetConst<{ SoundnessState::Unknown }>,
    proof: SoundnessProof,    // only wasm4pm produces this
) -> WfNetConst<{ SoundnessState::Witnessed }> { /* ... */ }
\end{lstlisting}

A caller who constructs a \texttt{WfNetConst<\{Unknown\}>} and passes it
directly to a function that demands \texttt{WfNetConst<\{Witnessed\}>} receives
a compile-time type error. No runtime assertion is needed; no boolean flag can
be forged.

\subsection{The Dual-Net Lesson}

During the AGI evaluation phase of the \texttt{wasm4pm-compat} development,
an evaluator proposed an alternative design with two separate types:
\texttt{WfNet} (an unparameterized workflow net) and \texttt{WfNetConst<S>}
(the const-parameterized version). The proposal argued that most callers would
use \texttt{WfNet} and only specialists would need \texttt{WfNetConst}.

The finding from evaluating this proposal was that a two-type design
reintroduces the verification gap. If \texttt{WfNet} and
\texttt{WfNetConst<S>} both exist, conversion functions between them are
required, and those conversion functions become the new forgery surface. A
caller who wishes to avoid the soundness gate converts their
\texttt{WfNetConst<\{Unknown\}>} to \texttt{WfNet} and passes the result to a
function that converts \texttt{WfNet} back to \texttt{WfNetConst<\{Witnessed\}>}
via an unsound stub. The type system cannot prevent this if the conversion
function exists.

The single-type design with a const-generic soundness parameter eliminates this
attack surface. There is exactly one type, \texttt{WfNetConst<SOUNDNESS>}, and
every function that requires witnessed soundness declares
\texttt{WfNetConst<\{SoundnessState::Witnessed\}>} in its signature. A caller
who cannot produce a \texttt{SoundnessProof} token cannot satisfy the bound.

% -----------------------------------------------------------------------------
\section{The LossPolicy Chain}
\label{sec:arch-losspolicy}
% -----------------------------------------------------------------------------

\subsection{Three Variants}

Some translations between process-evidence shapes cannot be lossless. The
canonical example is flattening an object-centric log (OCEL) to a classic
case-centric log (XES): a single object type must be chosen as the case notion,
and every event-to-object link to the other types is discarded. That discarded
structure is real evidence; it cannot vanish silently.

The \texttt{LossPolicy} enum forces the caller to declare, before any loss
occurs, how loss is to be handled:

\begin{lstlisting}[language=Rust, caption={\texttt{LossPolicy}: a pre-loss
  declaration. The default (via \texttt{Default::default()}) is
  \texttt{RefuseLoss} --- the safest option.}, label={lst:losspolicy}]
pub enum LossPolicy {
    /// Loss is not tolerated: refuse if evidence would be dropped.
    RefuseLoss,
    /// Loss is permitted under an explicitly named projection.
    AllowNamedProjection,
    /// Loss is permitted and must be itemised in a LossReport.
    AllowLossWithReport,
}

impl Default for LossPolicy {
    fn default() -> Self { LossPolicy::RefuseLoss }
}
\end{lstlisting}

The three variants form a lattice of increasing permissiveness.
\texttt{RefuseLoss} is the floor: no evidence may be dropped, and any
transformation that would drop evidence must instead return a refusal.
\texttt{AllowNamedProjection} names the projection but does not itemize the
loss; it is appropriate when the projection is well-understood and the loss
category is documented by the \texttt{ProjectionName}. \texttt{AllowLossWithReport}
requires the caller to produce a \texttt{LossReport} enumerating every
discarded item; it is the correct policy when the loss must be auditable.

\subsection{ProjectionName as a Semantic Newtype}

\texttt{ProjectionName} is a \texttt{\&'static str} newtype. It names a
projection in a stable, machine-readable way so that two runs of the same
named projection can be recognized as equivalent by downstream audit systems.
The use of \texttt{\&'static str} --- rather than \texttt{String} or a numeric
identifier --- ensures that names are baked into the binary, not computed from
runtime state that could vary between executions.

\begin{lstlisting}[language=Rust]
pub struct ProjectionName(pub &'static str);
// Example usage:
let name = ProjectionName("ocel-flatten-to-xes:by-order");
\end{lstlisting}

A \texttt{ProjectionName} must be chosen before the projection runs, not
after. This prevents a design where the name is derived from the result of the
projection, which would allow silent laundering of unnamed loss.

\subsection{LossReport Documentation}

A \texttt{LossReport<From, To, Items>} is the structured receipt of a lossy
projection. Its \texttt{From} and \texttt{To} type parameters tag the shapes
bridged (both are zero-sized \texttt{PhantomData} markers); \texttt{Items} is
the concrete record of discarded evidence (e.g., a \texttt{Vec<\&str>} of
dropped object-type names). The combination ensures that a report for an
OCEL-to-XES flattening cannot be confused with a report for a
process-tree-to-DFG projection.

\begin{lstlisting}[language=Rust, caption={A \texttt{LossReport} for
  an OCEL-to-XES flattening. The \texttt{From}/\texttt{To} phantom tags
  prevent the report from being confused with one for a different projection
  pair.}, label={lst:lossreport}]
enum OcelShape {}
enum XesShape  {}

let report = LossReport::<OcelShape, XesShape, Vec<&str>>::new(
    ProjectionName("ocel-flatten-to-xes:by-order"),
    LossPolicy::AllowLossWithReport,
    vec!["item", "invoice"],   // object types dropped
);

assert_eq!(report.policy, LossPolicy::AllowLossWithReport);
assert_eq!(report.lost,   vec!["item", "invoice"]);
assert!(!report.is_lossless());
\end{lstlisting}

\subsection{The Type-Level Guarantee That Silent Projection Is Impossible}

The \texttt{Project} trait is the only sanctioned path for a lossy
transformation. It requires:

\begin{enumerate}
  \item A \texttt{LossPolicy} argument --- decided by the caller before projection.
  \item A return type of \texttt{Result<LossReport<From, To, Lost>, Reason>}.
\end{enumerate}

There is no \texttt{Project} impl that returns a plain projected value without
either (a) a \texttt{LossReport} or (b) a named refusal. The constraint is
structural: any code that calls \texttt{project()} must handle both branches of
the \texttt{Result}. A caller who calls \texttt{project()} and discards the
\texttt{LossReport} by writing \texttt{let \_ = result.unwrap().lost;} has
explicitly acknowledged the loss --- the silence is named.

\begin{lstlisting}[language=Rust, caption={The \texttt{Project} trait.
  Silent loss is structurally impossible: every implementation must return
  either a \texttt{LossReport} or a named refusal.}, label={lst:project-trait}]
pub trait Project {
    type From;
    type To;
    type Lost;
    type Reason;

    fn project(
        self,
        policy: LossPolicy,
    ) -> Result<LossReport<Self::From, Self::To, Self::Lost>,
                Self::Reason>;
}
\end{lstlisting}

The diagnostic \texttt{CompatDiagnostic::LossyProjectionWithoutPolicy} and
\texttt{CompatDiagnostic::HiddenFlattening} are the named laws against which a
missing-policy violation and a silent-flattening violation are charged,
respectively.

\subsection{Algorithm: The Lawful Lossy Projection Protocol}

\begin{algorithm}[H]
  \caption{Lawful Lossy Projection}
  \label{alg:lossy-projection}
  \DontPrintSemicolon
  \KwIn{source evidence $e : \Evidence{T}{\texttt{Admitted}}{W}$,\quad
        policy $p : \texttt{LossPolicy}$,\quad
        projection name $n : \texttt{ProjectionName}$}
  \KwOut{projected evidence and loss report, or named refusal}
  \BlankLine
  $\text{dropped} \leftarrow \textsc{Identify\_Loss}(e)$\;
  \BlankLine
  \If{$p = \texttt{RefuseLoss}$ \textbf{and} $\text{dropped} \neq \emptyset$}{
    \Return $\Err(\textsc{Named\_Refusal\_Reason})$\;
  }
  \BlankLine
  $e' \leftarrow \textsc{Apply\_Projection}(e,\, n)$\;
  \BlankLine
  $r \leftarrow \texttt{LossReport::new}(n,\, p,\, \text{dropped})$\;
  \BlankLine
  \Return $\Ok\!\left(\, e'.\texttt{into\_projected}(),\ r \,\right)$\;
\end{algorithm}

The invariants enforced by the protocol are:
\begin{itemize}
  \item The loss is identified \emph{before} the projection is applied.
    A caller cannot observe the projected result and then decide retroactively
    whether the loss is acceptable.
  \item The \texttt{LossReport} is constructed from the \emph{same}
    \texttt{ProjectionName} that governs the projection. The name is fixed at
    call-time, not derived from the result.
  \item The refusal reason is a named type parameter (\texttt{Self::Reason}),
    not a string. This ensures that refusals are machine-readable and auditable.
\end{itemize}

\begin{remark}
  The protocol described in Algorithm~\ref{alg:lossy-projection} applies at
  the structural layer of \texttt{wasm4pm-compat}. It accounts for loss by
  shape; it does not run a discovery or conformance engine over the projected
  result. Engine-level actions on the projected evidence graduate to
  \texttt{wasm4pm}.
\end{remark}

% -----------------------------------------------------------------------------
\section{The Three-Layer Summary}
\label{sec:arch-summary}
% -----------------------------------------------------------------------------

Table~\ref{tab:three-layers} summarizes the three layers of the type system,
the mechanism each layer uses, and the invariant each layer enforces.

\begin{table}[h]
  \centering
  \caption{The three-layer type system of \texttt{wasm4pm-compat}.}
  \label{tab:three-layers}
  \begin{tabular}{@{}lll@{}}
    \toprule
    \textbf{Layer} & \textbf{Mechanism} & \textbf{Invariant enforced} \\
    \midrule
    State tokens
      & Empty enums as \texttt{PhantomData}
      & Cross-stage substitution is a type error \\
    Witness markers
      & Uninhabited enums + sealed \texttt{Witness} trait
      & Cross-authority substitution is a type error \\
    Const-generic law
      & \texttt{Require<\{EXPR\}>: IsTrue} where-bounds
      & Out-of-range quantitative invariants do not compile \\
    \bottomrule
  \end{tabular}
\end{table}

Together, the three layers produce a system in which every structural law of
the process-evidence domain --- lifecycle ordering, authority identity,
soundness non-forgeability, metric bounds, loss accountability --- is a
compile-time fact. No runtime state, no boolean flag, and no convention
agreed by comment alone can substitute for these invariants. The compiler
enforces them at every call site, without exception, at zero runtime cost.
